\paragraph*{04.03.03.05 Für das Projekt relevante fachliche Standards und Tools beurteilen, nutzen und weiterentwickeln}
\kcilabel{para:04.03.03.05_FachlicheStandardsundTools}{04.03.03.05}
\label{para:04.03.03.05_FachlicheStandardsundTools}


%%% Task
\par Als \gls{wsl} beurteilte, nutzte und entwickelte ich fachliche Standards und Tools weiter (SAP-Entwicklungs- und Transportstandards, Toolstrategie, \gls{bfa} für Schnittstellen, Ablösung SAP CALM zu JIRA/Confluence). Ziel: Konsistentes, effizientes Vorgehensmodell mit Qualität, Nachvollziehbarkeit, Transparenz und Effizienz.

%%% Action
\par Systematische Analyse und Bewertung nach Kriterien (Skalierbarkeit, Integrationsfähigkeit, Governance-Konformität, Nutzbarkeit). Migrationspfad für JIRA/Confluence, Governance-Regeln, Abstimmung mit Partnern. Etablierte Austauschformate (Show-and-Tell, Lessons-Learned, Workshops). Schulungen via \gls{myit} (vgl. \autoref{fig:ProgrammschulungMyIt}). Kombination aus Governance, Enablement und Verbesserungsprozess.

%%% Result
\par Steigerung von Qualität, Effizienz, Konsistenz; termingerechter Go-Live der \gls{ais}. Reduzierte Fehler durch Standards in \gls{is}, \gls{uac}, \gls{dp} (vgl. \autoref{fig:R2HBFA}). Höhere Transparenz und Rückverfolgbarkeit durch JIRA/Confluence. \gls{bfa} als tragfähig für Schnittstellenharmonisierung. Wissensaustausch professionalisierte Arbeitsweise und Übernahme in angrenzende Streams.

%%% Reflect
\par Frühzeitige Strategie entscheidend. Wechsel hätte durch frühere Stakeholder-Einbindung beschleunigt werden können.\endnote{Anders als SAP CALM war JIRA/Confluence bereits teilweise im \gls{adac} etabliert.} Kommunikationsstrategie hätte Akzeptanz von \gls{aif} gesteigert. Zukünftig: Frühe Toolentscheidungen, iterative Abstimmungen, kontinuierliche Weiterentwicklung als Governance-Prozess.